\subsection{History}

The idea of an international police body predates INTERPOL by nearly a decade.

In 1914, the first International Criminal Police Congress met in Monaco, where officials drafted early principles for cross-border cooperation and extradition,\source{deflem2002} but the First World War suspended the project. It revived in Vienna, where on 7 September 1923 delegates from twenty countries founded the International Criminal Police Commission (ICPC) under Johannes Schober, president of the Vienna police, in answer to a new generation of criminals who used cars, telephones, and open borders to outrun investigations that still halted at the frontier.\source{pace2017,interpolRepository2024} Through the interwar years, the Commission built the foundations of the modern system: a 1927 resolution urged each member to designate a central national contact point, the forerunner of today's NCB, and by the 1930s the organization had created specialized units and launched an international radio network.\source{interpolKeyDates}

The Second World War nearly destroyed it. After Germany annexed Austria in 1938, the Nazis seized control of the ICPC by deposing its president, and its leadership passed to a succession of senior SS officers; the headquarters were ultimately moved to Berlin, most member states withdrew, and the body was no longer functional.\source{deflem2002,interpolKeyDates}

\pullquote{This captured period is the reason the post-war organization treats political neutrality as foundational rather than decorative.}

In 1946, police officials meeting in Brussels reconstituted the body as the International Criminal Police Organization, with a new headquarters in Paris, and brought ``INTERPOL,'' long used as its telegraphic address, into common use.\source{deflem2002,britannica2024} The modern Constitution followed in 1956, formally establishing the name ICPO-INTERPOL\source[art.~1]{interpolConstitution} and the structure that governs the organization today, and the headquarters moved to Lyon in 1989.\source{britannica2024}

INTERPOL's membership and mandate have expanded steadily ever since. From twenty founding members in 1923, the organization grew to one hundred countries by 1967 and to its present 196 in 2023, making it second only to the United Nations in membership among international organizations.\source{wood2026} Its operational focus has widened as the character of crime kept changing, which brought us to this committee agenda of transnational financial crime.
